Partial RSA (Section~\ref{sec:gainmethods}) was used with covariate RDMs built from image statistics:
$|L_i - L_j|$ (mean luminance), $|C_i - C_j|$ (RMS contrast), the Euclidean distance between
coarse-luminance images (the grey image filtered to the lowest spatial-frequency band of
\citet{evalres}, including the mean, at $28 \times 28$; LP), the predecessor-luminance RDM
$|L_{\mathrm{prev}(i)} - L_{\mathrm{prev}(j)}|$, and the presentation-position RDM $|p_i - p_j|$
($T$). Ranks of the model and the brain RDM are regressed on the ranks of the covariates (with
intercept), and the residuals are correlated (Pearson).

The analyses of steps 44, 45, 47, 47b and 47c in THINGS: the decomposition of run centring
(Table~\ref{tab:decomp}), partial RSA with image-statistic, position and predecessor covariates
(Tables~\ref{tab:partial} and~\ref{tab:partialother}), and run-level quantities per subject
(Table~\ref{tab:runlevel}). The luminance correlation of the V1 RDM, divided by the leave-one-subject-out
lower bound, was \LumRatioTsev{} in THINGS-720, \LumRatioTest{} in the test images and \LumRatioNsd{} in NSD
(W4 met); for the coarse-luminance RDM the ratios were \LpRatioTsev, \LpRatioTest{} and \LpRatioNsd{} (W7
met). With run centring, the THINGS-720 ratio fell to \LumBoundRc{} (RC2 not met), and the multiverse had
\RuleRcFour{} distinct analyses with an interval above 0 (RC4 not met). The run-level variance of the
run-mean luminance was at the \RunLumPctSubOneTwo{}th percentile of random assignments for sub-01/02 and at
the \RunLumPctSubThree{}th for sub-03 (RC3 not met, reversed).

% generated by scripts/tables_v2.py; every number is a macro of numbers.tex
\begin{table}[H]
\centering
\small
\begin{adjustbox}{max width=\textwidth, max totalheight=0.82\textheight}
\begin{tabular}{@{}lrrrrr@{}}
\toprule
Variant (Section~\ref{sec:convmethods}) & sub-01 & sub-02 & sub-03 & mean V1 & mean LOC \\
\midrule
V0 nothing removed (paper cell) & $\DecVoneVZeroSubOne$ & $\DecVoneVZeroSubTwo$ & $\DecVoneVZeroSubThree$ & $\DecVoneVZero$ & $\DecLocVZero$ \\
V1 session-mean profile $m$ removed & $\DecVoneVOneSubOne$ & $\DecVoneVOneSubTwo$ & $\DecVoneVOneSubThree$ & $\DecVoneVOne$ & $\DecLocVOne$ \\
V2 run deviations $d_r$ removed & $\DecVoneVTwoSubOne$ & $\DecVoneVTwoSubTwo$ & $\DecVoneVTwoSubThree$ & $\DecVoneVTwo$ & $\DecLocVTwo$ \\
V3 both removed (run centring) & $\DecVoneVThreeSubOne$ & $\DecVoneVThreeSubTwo$ & $\DecVoneVThreeSubThree$ & $\DecVoneVThree$ & $\DecLocVThree$ \\
V4 = V3 + $d_{r\parallel}$ re-added &  &  &  & $\DecVoneVFour$ & $\DecLocVFour$ \\
V5 = V3 + $d_{r\perp}$ re-added &  &  &  & $\DecVoneVFive$ & $\DecLocVFive$ \\
\bottomrule
\end{tabular}
\end{adjustbox}
\caption{Decomposition of run centring in THINGS-720 (step 47c): $\Delta$ (untrained minus Backprop, paper cell, cross-run pairs) after removing parts of the voxel means. Per subject for V0--V3, V1 ROI.}
\label{tab:decomp}
\end{table}

% generated by scripts/tables_v2.py; every number is a macro of numbers.tex
\begin{table}[H]
\centering
\small
\begin{adjustbox}{max width=\textwidth, max totalheight=0.82\textheight}
\begin{tabular}{@{}lrr@{}}
\toprule
Covariates partialled (THINGS-720, paper cell) & V1 & LOC \\
\midrule
none & $\PartFfVoneNone$ & $\PartFfLocNone$ \\
$L$, $C$ (global luminance, RMS contrast) & $\PartFfVoneLc$ & $\PartFfLocLc$ \\
$T$ (presentation position) & $\PartFfVoneT$ & $\PartFfLocT$ \\
$L$, $C$, $T$ & $\PartFfVoneLct$ & $\PartFfLocLct$ \\
LP (coarse luminance image) & $\PartFffVoneLp$ & $\PartFffLocLp$ \\
$L$, $C$, LP & $\PartFffVoneLclp$ & $\PartFffLocLclp$ \\
$L$, $C$, LP, $T$ & $\PartFffVoneLclpt$ & $\PartFffLocLclpt$ \\
\midrule
run centring (step 47) & $\RunCentVoneRc$ & $\RunCentLocRc$ \\
run centring, then $L$, $C$ & $\RunCentVoneRcLc$ & $\RunCentLocRcLc$ \\
\bottomrule
\end{tabular}
\end{adjustbox}
\caption{Partial RSA in THINGS-720 (steps 44, 45, 47): $\Delta$ after rank-regressing both RDMs on covariate RDMs. THINGS test images and NSD: Table~\ref{tab:partialother}.}
\label{tab:partial}
\end{table}

% generated by scripts/tables_v2.py; every number is a macro of numbers.tex
\begin{table}[H]
\centering
\small
\begin{adjustbox}{max width=\textwidth, max totalheight=0.82\textheight}
\begin{tabular}{@{}lrrrr@{}}
\toprule
Covariates & test V1 ($k = 1$) & test LOC ($k = 1$) & NSD V1 ($k = 1$) & NSD LO-obj ($k = 1$) \\
\midrule
none & $\PartTestVoneNone$ & $\PartTestLocNone$ & $\PartNsdVoneNone$ & $\PartNsdLoobjNone$ \\
$L$, $C$ & $\PartTestVoneLc$ & $\PartTestLocLc$ & $\PartNsdVoneLc$ & $\PartNsdLoobjLc$ \\
LP & $\PartTestVoneLp$ & $\PartTestLocLp$ & $\PartNsdVoneLp$ & $\PartNsdLoobjLp$ \\
$L$, $C$, LP & $\PartTestVoneLclp$ & $\PartTestLocLclp$ & $\PartNsdVoneLclp$ & $\PartNsdLoobjLclp$ \\
\bottomrule
\end{tabular}
\end{adjustbox}
\caption{Partial RSA in the THINGS test images and in NSD (step 45), paper cell, standard normalisation of each data set.}
\label{tab:partialother}
\end{table}

% generated by scripts/tables_v2.py; every number is a macro of numbers.tex
\begin{table}[H]
\centering
\small
\begin{adjustbox}{max width=\textwidth, max totalheight=0.82\textheight}
\begin{tabular}{@{}lrrrrrrr@{}}
\toprule
Subject & $\Delta A$ & percentile & $\rho$(run RDM, run index) & run-level part of $\Delta$ & SD of $d_r\!\cdot\!\hat m$ & share along $\hat m$ & order $\rho$ \\
\midrule
sub-01 & $\MechDeltaASubOne$ & $\MechPctSubOne$ & $\MechRunIdxSubOne$ & $\MechRunPartSubOne$ & $\DecSdSubOne$ & $\DecShareSubOne\%$ & $\OrderRhoSubOne$ \\
sub-02 & $\MechDeltaASubTwo$ & $\MechPctSubTwo$ & $\MechRunIdxSubTwo$ & $\MechRunPartSubTwo$ & $\DecSdSubTwo$ & $\DecShareSubTwo\%$ & $\OrderRhoSubTwo$ \\
sub-03 & $\MechDeltaASubThree$ & $\MechPctSubThree$ & $\MechRunIdxSubThree$ & $\MechRunPartSubThree$ & $\DecSdSubThree$ & $\DecShareSubThree\%$ & $\OrderRhoSubThree$ \\
\bottomrule
\end{tabular}
\end{adjustbox}
\caption{Run-level quantities per subject (steps 44, 47b, 47c). $\Delta A$: alignment of the brain run RDM with the untrained minus the Backprop run RDM, and its percentile in the null of random run assignments; run-level part: $\Delta$ without minus with run centring; SD and share: the projection of the run deviations on the session-mean profile; order $\rho$: Spearman between the V1 RDM and the presentation-position RDM.}
\label{tab:runlevel}
\end{table}

